\SourceLeaf{008}
\begin{center}\small
A NEW PROOF OF THE THEOREM\\[5pt]
\textbf{THAT EVERY POLYNOMIAL}\\[5pt]
IN ONE VARIABLE\\[5pt]
CAN BE RESOLVED INTO REAL FACTORS OF THE FIRST OR SECOND DEGREE.\\[8pt]
\rule{1.5cm}{0.3pt}
\end{center}
\Article{1}
Every determinate algebraic equation can be reduced to the form
\[
x^m+Ax^{m-1}+Bx^{m-2}+\text{etc.}+M=0
\]
where $m$ is a positive integer. If we denote the left-hand side of this equation by $X$ and suppose that the equation $X=0$ is satisfied by several unequal values of $x$, say by setting $x=\alpha$, $x=\beta$, $x=\gamma$, etc., then the function $X$ will be divisible by the product of the factors $x-\alpha$, $x-\beta$, $x-\gamma$, etc. Conversely, if the product of several linear factors $x-\alpha$, $x-\beta$, $x-\gamma$, etc. divides the function $X$, then the equation $X=0$ will be satisfied by setting $x$ equal to any of the quantities $\alpha$, $\beta$, $\gamma$, etc. Finally, if $X$ equals a product of $m$ such linear factors (whether they are all different or some of them are identical), no linear factors other than these can divide the function $X$. An equation of degree $m$ therefore cannot have more than $m$ roots. At the same time, however, it is evident that an equation of degree $m$ can have \textit{fewer} roots even if $X$ can be resolved into $m$ linear factors: for if some of these factors are identical, the number of different ways of satisfying the equation will necessarily be smaller than $m$. Yet, for the sake of uniformity, mathematicians have preferred to say that the equation has $m$ roots in this case too, with only some of them being equal to one another; which they may certainly allow themselves to do.
